\documentclass[a4paper,12pt,titlepage]{article}

\usepackage[T1,T2A]{fontenc}     % форматы шрифтов
\usepackage[utf8]{inputenc}     % кодировка символов, используемая в данном файле
\usepackage[russian]{babel}      % пакет русификации
\usepackage{geometry}		 % для настройки размера полей
\usepackage{indentfirst}         % для отступа в первом абзаце секции
\usepackage{multirow}            % для таблицы с результатами
\usepackage{amsmath}
\usepackage{amssymb}
\usepackage{float}

% выбираем размер листа А4, все поля ставим по 3см
\geometry{a4paper,left=30mm,top=30mm,bottom=30mm,right=30mm}

\setcounter{secnumdepth}{0}      % отключаем нумерацию секций

\begin{document}

% Титульный лист
\begin{titlepage}
    \begin{center}
	{\small \sc Московский государственный университет \\имени М.~В.~Ломоносова\\
	Факультет вычислительной математики и кибернетики\\}
	\vfill
	{\Large \sc Отчет по заданию №1}\\
	~\\
	{\large \bf \guillemotleft{}Методы сортировки\guillemotright{}}\\ 
	~\\
	{\large \bf Вариант 8}
    \end{center}
    \begin{flushright}
	\vfill {Выполнил:\\
	Коцин~А.~О.}
    \end{flushright}
    \begin{center}
	\vfill
	{\small Москва\\2026}
    \end{center}
\end{titlepage}

% Автоматически генерируем оглавление на отдельной странице
\tableofcontents
\newpage

\section{Постановка задачи}

В данной работе требуется реализовать и экспериментально сравнить два алгоритма сортировки массивов целых чисел: пирамидальную сортировку (Heap Sort) и сортировку Шелла (Shell Sort). Оба алгоритма сортируют массив \textbf{по убыванию}.

В ходе экспериментов необходимо:

\begin{itemize}
\item реализовать оба метода сортировки с подсчётом количества операций сравнения и перемещения элементов;
\item провести тестирование на массивах различных размеров: $n = 10$, $100$, $1000$, $10000$;
\item использовать четыре типа входных данных для оценки поведения алгоритмов при различных степенях упорядоченности: отсортированные данные, данные в обратном порядке, два варианта случайных данных;
\item провести анализ полученных результатов и сравнить эффективность алгоритмов.
\end{itemize}

Сравнение производится по количеству операций сравнения и обмена элементов, что позволяет оценить трудоёмкость алгоритмов независимо от внешних факторов.

\newpage

\section{Результаты экспериментов}

В данном разделе приводятся результаты экспериментального сравнения Heap Sort и Shell Sort.

\subsection{Экспериментальные данные}

Эксперименты проводились на массивах размером $n = 10$, $100$, $1000$, $10000$ элементов, на четырёх типах входных данных:
\begin{itemize}
    \item \textbf{Sorted} — массив упорядочен по возрастанию
    \item \textbf{Reverse} — массив упорядочен по убыванию (обратный порядок)
    \item \textbf{Random1} — случайные данные (seed = 123)
    \item \textbf{Random2} — случайные данные (seed = 456)
\end{itemize}

Каждый тест повторялся 5 раз, результаты усредняются.

\subsubsection{Heap Sort}

\begin{table}[H]
\small
\centering
\begin{tabular}{|c|c|c|c|c|c|}
    \hline
    \multirow{2}{*}{\textbf{n}} & \multirow{2}{*}{\textbf{Параметр}} & \multicolumn{4}{|c|}{\textbf{Тип входных данных}} \\
    \cline{3-6}
    & & \textbf{Sorted} & \textbf{Reverse} & \textbf{Random1} & \textbf{Random2} \\
    \hline
    \multirow{2}{*}{10} & Сравнения & 35 & 41 & 37 & 39 \\
    \cline{2-6}
                        & Перемещения & 21 & 30 & 26 & 27 \\
    \hline
    \multirow{2}{*}{100} & Сравнения & 944 & 1081 & 1023 & 1038 \\
    \cline{2-6}
                         & Перемещения & 516 & 640 & 587 & 582 \\
    \hline
    \multirow{2}{*}{1000} & Сравнения & 15965 & 17583 & 16842 & 16856 \\
    \cline{2-6}
                          & Перемещения & 8316 & 9708 & 9090 & 9066 \\
    \hline
    \multirow{2}{*}{10000} & Сравнения & 226682 & 244460 & 235277 & 235329 \\
    \cline{2-6}
                           & Перемещения & 116696 & 131956 & 124088 & 124197 \\
    \hline
\end{tabular}
\caption{Результаты работы Heap Sort}
\end{table}

\subsubsection{Shell Sort}

\begin{table}[H]
\small
\centering
\begin{tabular}{|c|c|c|c|c|c|}
    \hline
    \multirow{2}{*}{\textbf{n}} & \multirow{2}{*}{\textbf{Параметр}} & \multicolumn{4}{|c|}{\textbf{Тип входных данных}} \\
    \cline{3-6}
    & & \textbf{Sorted} & \textbf{Reverse} & \textbf{Random1} & \textbf{Random2} \\
    \hline
    \multirow{2}{*}{10} & Сравнения & 27 & 22 & 29 & 29 \\
    \cline{2-6}
                        & Перемещения & 57 & 44 & 56 & 55 \\
    \hline
    \multirow{2}{*}{100} & Сравнения & 668 & 503 & 793 & 859 \\
    \cline{2-6}
                         & Перемещения & 1266 & 1006 & 1343 & 1416 \\
    \hline
    \multirow{2}{*}{1000} & Сравнения & 11716 & 8006 & 14866 & 15062 \\
    \cline{2-6}
                          & Перемещения & 20712 & 16012 & 23391 & 23560 \\
    \hline
    \multirow{2}{*}{10000} & Сравнения & 172578 & 120005 & 266889 & 262381 \\
    \cline{2-6}
                           & Перемещения & 302570 & 240010 & 391960 & 387497 \\
    \hline
\end{tabular}
\caption{Результаты работы Shell Sort}
\end{table}

\subsection{Анализ результатов}

На основе полученных экспериментальных данных можно сделать следующие наблюдения:

\textbf{Heap Sort:}
\begin{itemize}
    \item Демонстрирует стабильное поведение, почти независимое от типа входных данных.
    \item Количество сравнений и перемещений растёт примерно согласно $O(n \log n)$. Например, при увеличении $n$ от 1000 до 10000 (в 10 раз) количество сравнений увеличивается примерно в 14 раз, что близко к $10 \log_2 10 \approx 33$ с коэффициентом реализации.
    \item На отсортированных и обратно отсортированных данных показывает примерно одинаковые результаты.
\end{itemize}

\textbf{Shell Sort:}
\begin{itemize}
    \item Демонстрирует высокую адаптивность к типу входных данных.
    \item На отсортированных данных значительно быстрее, чем Heap Sort: 172578 сравнений против 226682 для Heap Sort при $n = 10000$.
    \item На обратно отсортированных данных также показывает улучшенные результаты: 120005 сравнений против 244460 для Heap Sort.
    \item На случайных данных показывает более высокое значение операций по сравнению с Heap Sort.
    \item Общая тенденция: меньше сравнений, но больше перемещений, что компенсирует выигрыш в сравнениях.
\end{itemize}



\section{Структура программы и спецификация функций}

Программа состоит из трёх основных файлов:

\subsection{sort.h}

Заголовочный файл, определяющий интерфейс модуля сортировки:

\begin{itemize}
    \item \texttt{void generate\_arr(int array[], int size, int type)} — генерирует массив заданного типа и размера
    \item \texttt{int* copy\_arr(const int src\_arr[], int size)} — создаёт копию массива
    \item \texttt{void heap\_sort(int array[], int size)} — сортирует массив методом кучи по убыванию
    \item \texttt{void shell\_sort(int array[], int size)} — сортирует массив методом Шелла по убыванию
    \item \texttt{void reset\_stats(void)} — обнуляет счётчики операций
    \item \texttt{long long get\_heap\_comps(void)} — возвращает количество сравнений в Heap Sort
    \item \texttt{long long get\_heap\_moves(void)} — возвращает количество перемещений в Heap Sort
    \item \texttt{long long get\_shell\_comps(void)} — возвращает количество сравнений в Shell Sort
    \item \texttt{long long get\_shell\_moves(void)} — возвращает количество перемещений в Shell Sort
\end{itemize}

\subsection{sort.c}

Реализация алгоритмов сортировки с подсчётом операций. Используются статические переменные для ведения счётчиков:

\begin{itemize}
    \item \texttt{static long long heap\_comps, heap\_moves} — счётчики для Heap Sort
    \item \texttt{static long long shell\_comps, shell\_moves} — счётчики для Shell Sort
\end{itemize}

\textbf{Heap Sort} реализована в две фазы:
\begin{enumerate}
    \item Построение минимальной кучи из исходного массива (функция \texttt{heapify})
    \item Итеративное извлечение минимума и перестройка кучи
\end{enumerate}

\textbf{Shell Sort} использует последовательность промежутков $n/2, n/4, \ldots, 1$ и для каждого промежутка выполняет сортировку вставками.

\subsection{main.c}

Программа тестирования, которая:

\begin{enumerate}
    \item Определяет размеры тестов: $n = 10, 100, 1000, 10000$
    \item Для каждого размера выполняет 5 итераций
    \item Для каждой итерации генерирует 4 типа массивов
    \item Вызывает оба алгоритма и собирает статистику
    \item Усредняет результаты и выводит таблицу
\end{enumerate}

\newpage

\section{Отладка программы, тестирование функций}

Для тестирования и отладки методов сортировки выполнены следующие проверки:

\begin{enumerate}
    \item \textbf{Тест 1: Малые массивы ($n = 10$)}
    \begin{itemize}
        \item Проверена корректность сортировки по убыванию
        \item Убедились, что счётчики инициализируются правильно
        \item Проверили работу на четырёх типах входных данных
    \end{itemize}
    
    \item \textbf{Тест 2: Средние массивы ($n = 100, 1000$)}
    \begin{itemize}
        \item Проверена масштабируемость алгоритмов
        \item Убедились, что результаты reproducible (повторяемые) при одинаковых seed'ах
        \item Проверена освобождение памяти (отсутствие memory leaks)
    \end{itemize}
    
    \item \textbf{Тест 3: Большие массивы ($n = 10000$)}
    \begin{itemize}
        \item Проверена корректность работы на больших объёмах данных
        \item Убедились в отсутствии buffer overflow'ов
        \item Проверена корректность вычисления среднего значения
    \end{itemize}
\end{enumerate}

Все тесты пройдены успешно. Программа работает стабильно без ошибок и утечек памяти.

\newpage

\section{Анализ допущенных ошибок}

При реализации программы критические ошибки не были обнаружены. Оба алгоритма — Heap Sort и Shell Sort — корректно реализуют соответствующие методы сортировки по убыванию. 

\bigskip
Программа успешно прошла тестирование на всех размерах входных массивов ($n = 10, 100, 1000, 10000$) и четырёх типах входных данных без потери корректности или утечек памяти.

\newpage

\section*{Выводы}

На основе проведённых экспериментов можно сделать следующие выводы:

\begin{enumerate}
    \item \textbf{Heap Sort} обеспечивает предсказуемую производительность $O(n \log n)$ независимо от типа входных данных, что делает его надёжным выбором при необходимости гарантированного времени выполнения.
    
    \item \textbf{Shell Sort} демонстрирует высокую адаптивность к входным данным, достигая значительного преимущества на отсортированных и частично упорядоченных массивах, что объясняется его базированием на сортировке вставками.
    
    \item Для практического применения рекомендуется:
    \begin{itemize}
        \item использовать Heap Sort, когда требуется гарантированная $O(n \log n)$ производительность;
        \item использовать Shell Sort для масивов среднего размера с вероятностью частичной упорядоченности;
        \item для очень больших наборов данных рассмотреть использование быстрой сортировки (Quick Sort) или сортировки слиянием (Merge Sort).
    \end{itemize}
    
    \item Подсчёт операций сравнения и перемещения — эффективный метод анализа алгоритмов, позволяющий объективно сравнивать их эффективность независимо от деталей реализации.
\end{enumerate}

\newpage
\begin{raggedright}
\addcontentsline{toc}{section}{Список цитируемой литературы}
\begin{thebibliography}{99}

\bibitem{heapsort} Heap Sort on GeeksforGeeks.
    \texttt{https://www.geeksforgeeks.org/heap-sort/}
    Дата обращения: 24 февраля 2026.

\bibitem{shellsort} Shell Sort on GeeksforGeeks.
    \texttt{https://www.geeksforgeeks.org/shellsort/}
    Дата обращения: 24 февраля 2026.

\bibitem{sorting} Comparison of Sorting Algorithms on Wikipedia.
    \texttt{https://en.wikipedia.org/wiki/Sorting\_algorithm}
    Дата обращения: 24 февраля 2026.

\bibitem{cs} Кормен Т., Лейзерсон Ч., Ривест Р., Штайн К. Алгоритмы: построение и анализ.
    Второе издание.~--- М.:<<Вильямс>>, 2005.

\end{thebibliography}
\end{raggedright}

\end{document}
